% Validation Methodology for Figure 1: Contextual Sensitivity (K=10)
% Target: KDD 2026

\subsection{Validation Methodology}
\label{sec:validation_methodology}

We validate the contextual routing signal through a regression-based
design with a distribution-free null baseline and a formal
likelihood-ratio test.  All analyses use the held-out test set
exclusively ($N{=}750$ prompts, $K{=}10$ models), with PCA trained on
independent data.

\subsubsection{Experimental Design}

\paragraph{Analysis pipeline matches the router.}
banditGPT's \texttt{FeatureService} encodes prompts with a sentence
transformer (\texttt{all-MiniLM-L6-v2}), compresses the 384D embeddings
to 32~dimensions via PCA, and appends a bias term for a 33-dimensional
feature vector.  Hybrid LinUCB then regresses these features against
per-model rewards to predict which model will perform best.  Our
validation tests the same relationship: does reward depend on PCA
features, and does the dependence differ across models?

\paragraph{Reward signal.}
Each prompt--model pair is scored by a panel of 3--4 GPT-4o judges.
Each judge emits a binary vote ($v \in \{0,1\}$) and a scalar confidence
($c \in [0,1]$).  The reward is $\text{mean}(v \times c)$ across judges,
yielding a continuous signal in $[0,1]$ that preserves evaluative nuance
lost by the binarised majority vote (see~\S\ref{sec:reward_signal} for
the full justification).

\paragraph{Regression specification.}
For each model $m$, we fit a simple OLS regression:
$r_{i,m} = \alpha_m + \gamma_m \, \widetilde{\text{PC1}}_i + \varepsilon_{i,m}$,
where $\widetilde{\text{PC1}}$ is the zero-mean, unit-variance first
principal component.  The slope $\gamma_m$ quantifies model~$m$'s
\emph{contextual sensitivity}: how much its reward changes per standard
deviation of PC1.  If all $\gamma_m$ are equal, context does not
differentiate models and a static policy suffices.

\paragraph{Likelihood-ratio test.}
To test whether slopes are heterogeneous, we compare a shared-slope null
$H_0$ against a per-model-slope alternative $H_1$ using $d{=}6$~PCA
components (Equations~\ref{eq:h0}--\ref{eq:h1}).  The test has
$(K{-}1) \times d = 54$ degrees of freedom.  Using six components rather
than just PC1 ensures the test captures multi-directional routing signal
(e.g., PC3 predicts Mixtral, PC31 predicts Llama-4-Maverick).

\paragraph{Null baseline.}
To rule out projection artifacts, we compute $|\rho|$ between oracle gain
and the first dimension of 100~independent random orthonormal projections
($384 \to 2$, QR-decomposed).  If the router PCA's $|\rho|$ falls within
this null distribution, the signal is artifactual; if it exceeds all
projections, the PCA concentrates genuine task-relevant variance.

\paragraph{Data independence.}
The PCA training data and the holdout evaluation data are entirely
independent:
\begin{itemize}[leftmargin=*,noitemsep,topsep=0pt]
    \item \textbf{PCA training:} 80K prompts from
          \texttt{routellm/gpt4\_judge\_battles}
          (HuggingFace)~\cite{ong2024routellm}.
    \item \textbf{Holdout evaluation:} 750 prompts scored by 10~models
          from the banditGPT evaluation pool, collected independently.
\end{itemize}
\noindent No decontamination step is needed---the datasets are disjoint
by provenance.

\subsubsection{Statistical Tests}

\paragraph{Primary tests.}
\begin{itemize}[leftmargin=*,noitemsep,topsep=0pt]
    \item \textbf{LR test} ($H_0$: shared slopes vs.\ $H_1$: per-model
          slopes; $d{=}6$ PCs, $\text{df}{=}54$): tests whether context
          affects models differentially.
    \item \textbf{Bootstrap CIs on $\gamma_m$} ($10\,000$ prompt-level
          resamples): quantifies per-model contextual sensitivity with
          uncertainty.  A slope is declared significant if its 95\% CI
          excludes zero.
\end{itemize}

\paragraph{Supporting tests.}
\begin{itemize}[leftmargin=*,noitemsep,topsep=0pt]
    \item \textbf{Spearman $\rho$} (PC1 vs.\ oracle gain): tests whether
          features predict where routing helps most.
    \item \textbf{Random projection null} ($N{=}100$): provides a
          distribution-free reference for the Spearman $|\rho|$.
\end{itemize}

\subsubsection{Design Rationale}

\paragraph{Why regression, not clustering.}
The bandit performs regression on continuous features.  A clustering-based
analysis introduces researcher degrees of freedom (threshold, cluster
count, distance metric) that are absent from the routing algorithm.
OLS regression tests the same functional form the bandit uses, making
the validation directly relevant.

\paragraph{Why PC1 in the visualisation.}
The LR test uses six PCA components, but Panel~A visualises PC1
alone for readability.  PC1 captures the most variance by construction
and yields the strongest single-component correlation for the plurality
of models.  Higher components contribute additional, complementary signal
(e.g., PC3, PC31) that strengthens the case for routing but would
overcomplicate a two-panel figure.

\paragraph{Multiple comparisons.}
The primary test is a single LR test (one $\chi^2$ statistic).
The ten per-model CIs in Panel~B are descriptive decompositions of the
omnibus result, not independent hypothesis tests; no multiplicity
correction is applied.  The Spearman correlation and random projection
null are confirmatory analyses supporting the LR test, not separate
hypotheses.

\subsubsection{Generalisability}

\paragraph{Holdout construction.}
The holdout ($N{=}750$) was created via stratified sampling across
category, complexity, and difficulty (Appendix~\ref{app:stratification}).
The difficulty dimension is derived from oracle reward statistics, which
enriches the sample for prompts with clear model differences.

\paragraph{Stratification sensitivity.}
The relative comparison between the router PCA and random projections is
unaffected by stratification: both are evaluated on the same holdout.
The absolute magnitude of $\gamma_m$ may be inflated relative to an
unstratified sample.
The existence of heterogeneous slopes is robust; the exact magnitudes
in deployment will depend on the user's prompt distribution.
